%% Proof of Secure Element Attestation
%% Build: pdflatex posea && pdflatex posea      (no bibtex run required)
\documentclass[conference,10pt]{IEEEtran}

\usepackage[T1]{fontenc}
\usepackage[utf8]{inputenc}
\usepackage{amsmath,amssymb}
\usepackage{booktabs}
\usepackage{microtype}
\usepackage{url}
\usepackage[hidelinks]{hyperref}
\usepackage{cite}

\emergencystretch=1.2em
\hyphenpenalty=1000

\newtheorem{observation}{Observation}
\newtheorem{definition}{Definition}
\newtheorem{claim}{Claim}

\begin{document}

\title{Proof of Secure Element Attestation:\\
Sybil Resistance from Hardware Scarcity Rather Than Cost}

\author{\IEEEauthorblockN{Jan Ku\v{c}era}
\IEEEauthorblockA{\textit{NADO}\\
\texttt{admin@bismuth.cz}}}

\maketitle

\begin{abstract}
Permissionless systems that reward participants rather than capital face an adversary ordinary
countermeasures cannot price: the identity farm. Every per-identity admission rule---registration
work, waiting periods, address caps---costs an adversary holding $n$ identities exactly
what it costs an honest participant holding one, per identity, and the adversary amortises fixed
setup while the honest participant cannot. In one deployment, some $1{,}000$ of
$1{,}191$ identities belonged to two or three operators running browser farms, capturing $42\%$
of emission.

Proof of Secure Element Attestation (PoSEA) makes identities \emph{scarce} rather than expensive,
admitting a participant on a hardware attestation verified deterministically and offline against
pinned vendor roots, one identity per physical device. We then address its largest gap: of $664$
measured attempts, $26.8\%$ came from machines with a healthy TPM~2.0 that no platform will
certify, the vendor's attestation-key service having no authority registered for the chip. Their
vendor-signed \emph{endorsement} key cannot sign, which is why the standard remedy is a privacy
certificate authority holding a key able to assert any chip's identity.

We give an enrolment protocol admitting such machines with no certificate authority and no
long-lived key at all, resting on the observation that \texttt{TPM2\_MakeCredential}'s credential
blob is deterministic in its (seed, name, secret) triple: a challenge issued once is replayable
offline by every verifier forever, so a commit--reveal across $k$ beacon-drawn challengers,
ordered by the ledger, supplies the freshness a signature would otherwise provide. Its security
reduces to publication order and to one TPM object attribute. We validate it against a TPM~2.0
simulator and physical silicon, and state the cost: automatic enrolment removes the human tap, so
identity price collapses onto the price of a chip.
\end{abstract}

\begin{IEEEkeywords}
Sybil resistance, hardware attestation, trusted platform module, permissionless consensus,
commit--reveal
\end{IEEEkeywords}

\section{Introduction}

A permissionless network that wishes to reward people rather than capital must decide who
counts as a person. Proof of Work~\cite{nakamoto} answers by making participation expensive in
energy; Proof of Stake answers by making it expensive in capital. Both are auctions, and in an
auction the participant holding the most of the scarce input wins. Concentration follows by
construction rather than by accident.

Systems that attempt to escape this by paying per participant inherit the Sybil
problem~\cite{douceur} in its sharpest form: if an identity can be minted by a script, a reward
paid per identity is a subsidy to whoever mints the most. This is not hypothetical. NANO's
representative spam, Idena's validation ceremonies~\cite{idena} and Nyzo's cycle each
encountered it, and each responded with additional per-identity rules that the next operator
circumvented by the method that defeated the last. We report a further instance, with
measurements, and argue that the recurrence is structural rather than a sequence of
implementation errors.

\subsection{Contributions}

\begin{itemize}
  \item We state and measure the linearity argument that makes per-identity admission rules
        ineffective against identity farms (\S\ref{sec:linearity}), with deployment data from a
        live chain on which four successive such rules failed.
  \item We specify PoSEA, an admission predicate based on secure-element attestation with
        deterministic, offline, replicable verification, and the per-device binding that makes
        ``one device, one identity'' enforceable (\S\ref{sec:posea}).
  \item We quantify the mechanism's coverage gap: $26.8\%$ of admission attempts are machines
        with functioning TPM~2.0 hardware that no platform will certify
        (\S\ref{sec:gap}).
  \item We give a certificate-authority-free enrolment protocol closing that gap
        (\S\ref{sec:noca}). Its novelty is not cryptographic: it is the observation that a TPM
        credential blob is deterministic in its inputs, which converts an inherently
        interactive possession proof into an artifact any verifier can re-derive, so that
        ledger \emph{ordering} can replace an issuer's signature.
  \item We analyse its security (\S\ref{sec:security}), reporting that it reduces to
        publication order and to the \texttt{restricted} object attribute, and validate it
        against a TPM~2.0 simulator and physical silicon (\S\ref{sec:eval}).
  \item We state the costs without mitigation (\S\ref{sec:discussion}), including one the
        CA-free path introduces: it removes the human tap, so the security argument shifts from
        non-automatable human work to hardware price alone.
\end{itemize}

\section{Background and threat model}

\subsection{Deployment context}

NADO is a permissionless chain with a two-lane block production schedule. A \emph{bonded} lane
weights producers by stake. An \emph{open} lane, receiving $30\%$ of slots, was designed so a
participant with no capital could produce blocks from an ordinary mobile device. Selection in
both lanes is a deterministic beacon-keyed weighted draw: one hash determines each slot's
producer, so there is no grinding advantage and no hash race. Finality is a bonded-stake quorum
with Casper-style epoch checkpoints~\cite{casper}.

The open lane was defended, successively, by a non-parallelisable sequential work proof at
registration, a per-subnet ingress budget, a per-address identity cap, and a probationary
period before a new identity could earn. Each was deployed in response to observed farming.
Each failed.

\subsection{Why per-identity rules fail}
\label{sec:linearity}

\begin{observation}[Linearity]
\label{obs:linearity}
Let $c$ be the cost an admission rule imposes per identity and $r$ the reward per identity. An
adversary controlling $n$ identities pays $nc$ and earns $nr$; an honest participant pays $c$
and earns $r$. The cost-to-reward ratio is $c/r$ for both. Raising $c$ therefore does not
disadvantage the adversary relative to the honest participant; it disadvantages both equally,
and prices out the honest participant first, because the adversary amortises fixed setup across
$n$ while the honest participant cannot.
\end{observation}

Observation~\ref{obs:linearity} is elementary. We state it explicitly because practice
continues to propose per-identity rules as anti-Sybil measures, and because it identifies what
a working defence must do: bound $n$ itself, rather than tax it.

\subsection{Measurement}

On 2026-09-06 the open lane carried $1{,}191$ registered identities. Approximately $1{,}000$
were attributable to two or three operators running headless browser farms on rented Linux and
Windows servers. The farmed set captured approximately $42\%$ of total emission. Every honest
participant was, in effect, subsidising three people.

\subsection{Threat model}
\label{sec:threat}

The adversary is computationally bounded, may hold arbitrary capital, may rent unlimited
conventional compute, and controls every host it operates, including their operating systems.
It may create unlimited network identities and key pairs at negligible cost, may observe all
ledger state, and may participate in the protocol in any role including as a validator.

We assume the adversary cannot (i) forge a signature from a silicon vendor's attestation root,
(ii) extract a private key from a secure element it does not physically possess, or (iii) cause
a genuine secure element to attest to a key that does not reside in it. Assumption (ii) is
bounded, not absolute, and \S\ref{sec:open} treats it as the construction's principal open
question; published attacks require physical access~\cite{faultpm,tpmfail}, which prices
extraction per device rather than granting a remote oracle---so far.

We do not assume any honest majority among challengers in \S\ref{sec:noca}; we assume only that
not all $k$ drawn challengers collude, and \S\ref{sec:security} makes the dependence explicit.
We do not assume verifier clocks agree: all validity is evaluated at a ledger-agreed anchor
time.

\section{PoSEA}
\label{sec:posea}

\subsection{Intuition}

Only two properties resist replication by an adversary with money but no people: capital, which
the bonded lane already weighs, and a physically distinct object. PoSEA anchors open-lane
identity in the second. A secure element is a physically distinct object that can attest to its
own genuineness; it cannot be instantiated in software, and---the property that matters
here---it cannot be instantiated in software by an adversary who already controls the host.

\begin{definition}[PoSEA admission]
\label{def:admission}
An identity is admitted for one lease period if it presents a WebAuthn~\cite{webauthn}
attestation statement in which
\begin{enumerate}
  \item a non-exportable key was generated inside a secure element;
  \item that key signs a challenge derived from a ledger-selected beacon value, so the
        signature cannot be precomputed;
  \item the attestation certificate chain terminates at a vendor root pinned as a protocol
        constant, every signature verifying at the anchor time;
  \item the attested boot state indicates a locked bootloader and verified
        boot~\cite{androidkeyattest}, where the class reports it; and
  \item where the class exposes a per-device certificate, that certificate is not currently
        bound to another identity.
\end{enumerate}
\end{definition}

\subsection{Device classes}

Four classes expose hardware attestation to a web page; they differ in whether they yield a
per-device handle, which determines whether condition~(5) is enforceable at all.

\begin{table}[t]
\caption{Device classes and per-device binding}
\label{tab:classes}
\centering
\small
\begin{tabular}{@{}llc@{}}
\toprule
Class & \texttt{fmt} & Bindable by \\
\midrule
Android 12+, locked & \texttt{android-key} & RKP cert \\
Windows, physical TPM & \texttt{tpm} & AIK cert \\
Any TPM 2.0, no CA (\S\ref{sec:noca}) & --- & endorsement key \\
FIDO2 security key & \texttt{packed} & \emph{not bindable} \\
Apple Secure Enclave & \texttt{apple} & \emph{no attestation} \\
\bottomrule
\end{tabular}
\end{table}

With remote key provisioning (Android 12+), the certificate at \texttt{x5c[1]} is unique to the
device and reused for every credential it creates until rotation---approximately two weeks in
our sample. A pre-RKP device instead carries a multi-year intermediate shared by up to $10^5$
units; the two are distinguished by validity span, and the batch case is refused, since binding
a shared certificate to one identity would lock out every other owner of that model. On
Windows, the attestation identity key (AIK) certificate is unique to the pair (TPM, account).

FIDO2 authenticators issue one batch certificate per $\geq 10^5$ units by design, as a privacy
measure: they identify a model, not a unit, and are refused. Apple passkeys return
\texttt{fmt: none} and carry no attestation through a web page (iOS 16+, macOS 13+, confirmed
2026-09-07); a native App~Attest bridge is the only route and is not deployed. What cannot be
bound is not evidence of anything, so both are excluded rather than admitted with a weaker
guarantee.

\subsection{Determinism as a requirement, not a preference}
\label{sec:determinism}

In a replicated setting, admission is consensus input: every verifier must reach the same
verdict on the same bytes, forever, including on replay years later. This forces four
decisions that a single-verifier deployment could take differently.

Roots are pinned by SHA-256 over DER and stored with the protocol. They are never fetched,
never learned from a peer, and change only by a gated commit. Revocation lists are deliberately
not consulted, because doing so makes validation network-dependent and time-dependent; a
compromised intermediate is handled as a root-set change. This trades revocation latency for
deterministic replay, consciously.

Certificate validity is evaluated at a ledger-agreed anchor time rather than a wall clock. With
wall clocks, a certificate expiring mid-block is valid on one verifier and expired on the next:
a fork with no adversary involved.

The attestation parser must reject duplicate CBOR map keys. Permitting them allows a statement
to verify as one chain and bind as another. This, and the anchor-time rule, correspond to holes
found and closed in the deployment.

Finally, a binding must occupy its per-device uniqueness key \emph{within} the block that
creates it, so two provers cannot bind the same device in one block by each checking only
against parent state.

\subsection{What the mechanism does and does not prove}

It proves the registration was produced on genuine hardware of a stated class, by a key
resident in that hardware, over a challenge that could not be precomputed; and on renewal, that
the same device key was present later.

It does not prove personhood. One human with one genuine device may create as many identities
as they physically tap, because WebAuthn requires user activation per signature. The claim is
that an identity costs \emph{non-automatable human work on genuine hardware}, not that
identities map one-to-one onto people; proof-of-personhood in the stronger sense remains a
separate problem~\cite{pop,offline-pseudonyms}. Rooting is not a loophole, and is worth stating
because it is the first objection raised: a rooted device is precisely what would automate the
tap, and a rooted device cannot attest, because key attestation reports boot state from the
secure element itself.

Block production is unchanged. Attestation decides membership of the draw, not its outcome, so
faster hardware confers no advantage; attestation and presence weight are excluded from
fork-choice weight and from the finality quorum. Presence earns reward and never purchases a
say in what is final.

\section{The uncertified-chip gap}
\label{sec:gap}

Table~\ref{tab:gap} gives the outcome of $664$ consecutive admission attempts.

\begin{table}[t]
\caption{Admission attempts by outcome ($n=664$)}
\label{tab:gap}
\centering
\small
\begin{tabular}{@{}lr@{}}
\toprule
Outcome & Share \\
\midrule
Succeeds (phone / hardware wallet) & $47.9\%$ \\
\textbf{TPM healthy, platform will not attest} & $\mathbf{26.8\%}$ \\
User-fixable (password manager took the ceremony) & $18.8\%$ \\
Succeeds (Windows Hello attests) & $6.5\%$ \\
\bottomrule
\end{tabular}
\end{table}

The second row is not a hardware fault. A WebAuthn \texttt{tpm} statement requires
\texttt{certInfo} to be signed by the key in \texttt{x5c[0]}, so the chip needs an
attestation-key \emph{certificate}, which on Windows means the platform vendor's attestation
service. For a large population of otherwise healthy machines, that service has no certificate
authority registered for the chip's KeyId and answers HTTP~404. The vendor has confirmed the
condition as service-side, with KeyIds not registrable by end users and no supported route for
a user to alter the trust database. No client-side change repairs it. On Linux the situation is
simpler and worse: there is no equivalent service and there has never been a path.

We established that the platform path could not be repaired from the client side by decoding
the platform's own key-attestation claim structure, which we document
separately~\cite{claimformat}: the platform API mints a fresh per-call signer whose certificate
no relying party has ever seen, so no arrangement of client-side calls can satisfy WebAuthn's
requirement.

These machines are not short of evidence. They carry a vendor-signed \emph{endorsement}
certificate---AMD's chaining to \texttt{CN=AMDTPM}, Intel's through the CSME issuing
authorities the chip holds in its own non-volatile memory. The hardware proof exists and is
verifiable offline against pinned roots. What is missing is one statement: \emph{this
attestation key lives in that certified chip}.

\section{Enrolment without a certificate authority}
\label{sec:noca}

\subsection{Why the obvious remedies fail}

Two proposals recur and both must be rejected.

\emph{Present the endorsement certificate as the proof.} It fails twice: the certificate is
public data, so anyone who has seen a machine's certificate can present it; and the endorsement
key \emph{cannot sign}. TCG defines it as a restricted decryption key~\cite{tcgek}, and real
certificates carry \texttt{Key Usage: critical, Key Encipherment} with \texttt{digitalSignature}
absent. It can never occupy the \texttt{x5c[0]} slot.

\emph{Derive the challenge from public ledger data,} so no round trip is needed. Anything every
verifier can recompute, the prover can recompute. The resulting proof proves nothing.

Proving possession of a decryption-only key has exactly one shape: send it something only it can
decrypt and observe it return. That is \texttt{TPM2\_ActivateCredential}~\cite{tpm2}, and it is
inherently interactive---which is why TCG specified the attestation-key-plus-privacy-CA model,
and why the failing service exists at all.

\subsection{Replacing the authority with an ordering}

A certificate authority converts the interactive proof into a durable artifact by signing it.
That works, and it is expensive in a sense unrelated to computation: a long-lived key able to
assert any endorsement identity is a key able to mint identities, and must be guarded as such
for the lifetime of the system. In a permissionless setting it is also a single party the whole
mechanism defers to, which is the objection the mechanism exists to answer.

\begin{observation}[Replayable challenges]
\label{obs:det}
The credential blob returned by \texttt{TPM2\_MakeCredential} is a deterministic function of the
triple (seed, object name, secret). Only the OAEP-wrapped seed is randomised, and no verifier
requires that half.
\end{observation}

Observation~\ref{obs:det} follows directly from the credential-protection construction of TPM
2.0 Part~1~\S24~\cite{tpm2}: the blob is a symmetric encryption of the secret under a key
derived from (seed, name) by KDFa, with an HMAC over the ciphertext and the name under a second
derivation from the same pair. Every input is fixed by the triple. We confirmed it empirically:
two calls with the same seed produce identical blobs and different wrapped seeds.

It is what removes the authority. A challenge issued once can be recomputed by every verifier
afterwards, offline and indefinitely, from values that are public after the fact. The
interactive proof therefore need not be signed into an artifact; it need only be
\emph{ordered}---and a ledger already orders things.

\subsection{The protocol}

\begin{definition}[CA-free enrolment]
\label{def:enrol}
Let $E$ be a chip's endorsement key, certified by a silicon vendor, and $A$ an attestation key
in the same chip with object name $N$. An enrolment is the following sequence of published
messages, each strictly later in the ledger order than the message it answers:
\begin{enumerate}
  \item the prover publishes $E$'s certificate chain and $A$'s public area, from which $N$ is
        derived;
  \item each of $k$ drawn challengers publishes
        $b_i = \texttt{MakeCredential}(E, N, S_i; R_i)$ together with the wrapped seed;
  \item the prover publishes $H(S_1 \Vert \cdots \Vert S_k)$, having recovered every $S_i$ by
        \texttt{TPM2\_ActivateCredential};
  \item each challenger publishes $(S_i, R_i)$.
\end{enumerate}
Verifiers accept when every $b_i$ is reproduced by $(S_i, R_i)$ under $(E, N)$ and the
commitment matches the concatenation in canonical challenger order.
\end{definition}

Nothing in Definition~\ref{def:enrol} is signed and no key persists between enrolments. There
is nothing to steal, to rotate, or to guard.

\subsection{Challenger selection}
\label{sec:draw}

The residual attack is a challenger privately leaking $S_i$ to a prover holding no chip. The
defence is not to trust challengers, but to require several and to draw rather than accept
them: the prover must recover every $S_i$, so forgery requires all $k$ to collude. This
replaces a key someone promises to protect with a property the system can observe.

The draw is weighted by a quantity the adversary cannot manufacture cheaply---in the
deployment, bonded stake---and keyed on the same grind-resistant beacon that keys producer
selection, so it is unpredictable before the enrolment is published and identical for every
verifier. An unweighted draw over a permissionless set is the obvious error: cheap identities
crowd the ballot and an adversary able to mint them draws its own challengers, restoring the
attack in full. The draw is without replacement, since seating one challenger twice halves the
collusion forgery requires. A short set is refused rather than accepted: an adversary able to
shrink the challenger pool must not thereby reduce what forgery costs.

\subsection{Binding and freshness}

Identity binds to the \emph{endorsement} key, never to an attestation key. A chip holds
unlimited attestation keys, so binding to one would hand a single machine an identity per
enrolment; the handle is SHA-256 over the endorsement certificate's
\texttt{SubjectPublicKeyInfo}. Regenerating the endorsement seed to fake a new chip changes
that key, and the vendor's certificate no longer matches it, so the fake cannot enrol at all.

An enrolment records that a key lives in a certified chip; it says nothing about when. Each
admission therefore presents a fresh \texttt{TPM2\_Certify} under the enrolled key whose
\texttt{extraData} carries that admission's own challenge. The verifier additionally checks the
\texttt{TPM\_GENERATED} magic, the \texttt{TPM\_ST\_ATTEST\_CERTIFY} structure type---a
different type places different fields at the same offsets, so accepting one would mean parsing
some other structure as this one---and that the certified object is the enrolled key rather
than another object in the chip.

\section{Security analysis}
\label{sec:security}

\subsection{The ordering is the proof}

\begin{claim}
\label{claim:order}
The equations of Definition~\ref{def:enrol}, verified without regard to publication order, are
satisfiable by a prover holding no TPM.
\end{claim}

Such a prover chooses $S_i$ and $R_i$ itself and computes $b_i$ from the endorsement public key,
which is public. Every check passes: each $b_i$ is reproduced by its $(S_i, R_i)$, and the
commitment matches by construction. Presented with all four messages simultaneously, a verifier
cannot distinguish this from a genuine enrolment.

What separates them is that message~3 was published after every message~2 and before every
message~4. The prover cannot learn $S_i$ except by holding the chip, and must commit before the
reveal, so it cannot read the answer off the ledger. Message~2 cannot precede message~1,
because \texttt{MakeCredential} seals to a specific object name, which is a digest of the
public area published in message~1. The four messages are therefore minimal: no adjacent pair
may be merged.

Two consequences are easy to implement incorrectly. First, equal ordering positions are not
ordered: two messages in one block have no order verifiers agree on, so every comparison must
be strict. Second, the commitment must be over all $k$ secrets jointly. One commitment per
secret would let a prover answer whichever challenges it managed to open and abandon the rest,
which is the $k$-of-$k$ requirement dissolving into $1$-of-$k$.

\subsection{The attribute the argument rests on}

\begin{claim}
\label{claim:restricted}
If an unrestricted signing key may be enrolled, the certify of \S\ref{sec:noca} carries no
information about where a key resides.
\end{claim}

A \texttt{restricted} signing key signs only structures the TPM itself generated. That is what
makes a certify evidence rather than merely a signature. An unrestricted key signs whatever its
host composes, including a forged \texttt{TPMS\_ATTEST} certifying a key that never existed in
any chip; a verifier would then accept the forgery under a genuinely enrolled name. The public
area is therefore validated before any challenge is issued, and must further carry
\texttt{sensitiveDataOrigin}---absent it, the private half may have been generated outside and
imported---and \texttt{fixedTPM} with \texttt{fixedParent}, absent which the key is duplicable
to another chip and one enrolment would licence every machine it is copied to.

\subsection{Attacks considered}

Table~\ref{tab:attacks} lists the attacks the construction refuses and the mechanism that
refuses each. We include the case it does \emph{not} refuse.

\begin{table}[t]
\caption{Attacks and the mechanism that refuses them}
\label{tab:attacks}
\centering
\small
\begin{tabular}{@{}p{0.44\columnwidth}p{0.46\columnwidth}@{}}
\toprule
Attack & Refused by \\
\midrule
Replay a copied endorsement certificate & Possession proof: the seed unwraps only under the endorsement private key \\
Prover seals its own challenge & Challenger set is drawn, not chosen (\S\ref{sec:draw}) \\
Prover commits after seeing a reveal & Strict ordering, message 3 before 4 \\
Prover opens some challenges, guesses the rest & Joint commitment over all $k$ secrets \\
Challenger reveals a different secret & Blob already published; MakeCredential is deterministic \\
Enrol a second key for a second identity & Binding is on the endorsement key \\
Reuse an old certify & \texttt{extraData} carries the admission's own challenge \\
Certify some other object in the chip & Certified name must equal the enrolled name \\
Enrol a key that can sign arbitrary data & \texttt{restricted} required (Claim~\ref{claim:restricted}) \\
Duplicate an enrolled key to another chip & \texttt{fixedTPM}, \texttt{fixedParent} required \\
\midrule
\textbf{All $k$ challengers collude} & \textbf{Not refused} (\S\ref{sec:security}) \\
\textbf{Physical key extraction} & \textbf{Not refused} (\S\ref{sec:open}) \\
\bottomrule
\end{tabular}
\end{table}

Collusion of all $k$ drawn challengers admits a prover with no chip. We do not claim otherwise.
The design intent is that this is a condition the system can \emph{observe}---the drawn set is
public, weighted and beacon-keyed---rather than a key whose custody someone asserts. With
$k=3$, the smallest value at which no single challenger can admit identities alone, an
adversary must both control challengers and win the draw for the specific enrolment it wishes
to forge.

\section{Implementation}
\label{sec:impl}

The reference implementation is dependency-free Python: KDFa, AES-CFB, MGF1, RSA-OAEP, PKCS\#1
v1.5 verification and the DER encoding are written out rather than imported. This is not
minimalism for its own sake. A verifier that must agree byte-for-byte with every other verifier
is poorly served by whichever version of a library each host happens to have installed, and in
the deployment this code runs inside consensus on nodes whose runtime carries no cryptographic
library at all. We record two obstacles that cost us time and that we expect to cost others the
same.

\subsection{Real endorsement certificates are not strictly DER}

AMD encodes \texttt{critical: FALSE} explicitly where DER requires the default omitted. Python's
\texttt{cryptography} refuses such a certificate outright with an \texttt{EncodedDefault} parse
error, while \texttt{openssl} and Rust's \texttt{x509-parser} accept it. An implementer
reaching for the obvious library will conclude the hardware is defective; it is not, and the
certificate is the one the vendor issued. We moved endorsement-chain parsing to a lenient
parser, shared by all verifiers.

\subsection{A published ``root'' that is an intermediate}

The certificate one vendor publishes as its endorsement root is cross-signed by a commercial
certificate authority and is therefore an intermediate. Pinning it would silently relocate the
trust decision to that authority, which is a different assertion from ``this vendor
manufactured this chip''. We verify self-signedness before pinning, and omit that vendor rather
than pin an intermediate.

\subsection{Signature verification}

PKCS\#1 v1.5 verification compares the entire re-encoded block~\cite{rfc8017}. An
implementation that instead searches for the digest within the padding accepts trivial
forgeries against small exponents---Bleichenbacher's 2006 result~\cite{finney,cve2006}, which
production libraries have shipped more than once. We pin this as a test rather than a comment.

\section{Evaluation}
\label{sec:eval}

We exercised the construction against a TPM~2.0 simulator, driving our own command encodings
over the simulator's transport. The TCG endorsement template derives its $314$-byte public
area; a restricted RSA-2048 signing key derives; our challenger's credential blob is accepted
and opened by \texttt{ActivateCredential}, returning the sealed secret; an independent verifier
re-derives the identical blob from the revealed pair; and a certify verifies, while a certify
produced by a \emph{different} key in the same chip is refused under the enrolled key's public
area. The full four-message enrolment reaches its proven state on values the chip actually
produced. The enrolment has additionally been completed against a physical AMD firmware TPM.

Simulation does not establish that a particular vendor's silicon behaves this way: a simulator
accepts templates a firmware TPM may reject, and carries no endorsement certificate at all, so
the one thing this path ultimately rests on---the vendor signature---cannot be exercised there.
This is why we report the physical result separately and regard broader hardware coverage as
required work rather than a formality.

\section{Discussion}
\label{sec:discussion}

\subsection{Vendor authority in the trust path}

Verification requires trusting that vendor attestation roots sign only genuine hardware. This
is a centralisation cost in a system whose premise is that no such authority should be
necessary, and we do not think it can be argued away.

We observe only that every Sybil mechanism trusts something: Proof of Work trusts that no
coalition controls a hash majority, an assumption weakening empirically; Proof of Stake trusts
that concentrated holders do not collude. PoSEA trusts pinned vendor roots, verified offline,
with no service able to revoke a participant at validation time. We regard this as a different
trust assumption rather than a strictly worse one, and note that reasonable readers disagree.

\subsection{Exclusion of modified devices}

A device whose owner has unlocked its bootloader cannot attest. This excludes,
disproportionately, the users most committed to the principle that an owner should control
their own hardware---the constituency most sympathetic to a permissionless chain. The
deployment mitigates this only partially, by a bonded lane requiring no device, which is a
capital gate and therefore not a substitute.

\subsection{What the CA-free path costs}

The human tap is gone. Enrolment and renewal are automatic, which is required for headless
machines and is why a server can participate at all; it also means a farm's chips cost nothing
to keep running. Identity \emph{count} is unchanged---one chip, one identity, however many
attestation keys it enrols---but the security argument shifts. Where the WebAuthn path prices
an identity at non-automatable human work on genuine hardware, this path prices it at the
hardware alone.

The choice of pinned endorsement roots therefore becomes the entire dial, and it reverses only
in the widening direction: narrowing the set later strands everyone already enrolled under the
wider one. A discrete TPM module costs roughly an order of magnitude less than the machine a
CPU-vendor firmware TPM implies. We regard this as a policy choice a deployment must make
deliberately and state, not a parameter to be set by pinning every root that can be found.

\section{Open problem}
\label{sec:open}

The security of the construction reduces to an empirical quantity we cannot presently bound:
the cost of extracting an attestation or endorsement private key from a secure element, and
whether that cost is specific to a model or generalises across a chip family.

The question is not open in the sense of untouched. TPM-FAIL~\cite{tpmfail} recovered ECDSA
private keys from an Intel firmware TPM and an STMicroelectronics TPM by timing analysis;
faulTPM~\cite{faultpm} extracted secrets from AMD firmware TPMs by voltage fault injection
against the AMD Secure Processor; and Samsung's TrustZone keymaster was shown to permit
attestation-relevant key compromise by design error~\cite{trustzone}. These results matter to
us directly: one of them targets a vendor whose root we pin.

What they establish is that extraction is possible with \emph{physical access to the target
machine}, which prices it per device and does not hand a remote adversary an oracle. What they
do not establish---and what we need---is whether any of these generalises to a key that can be
produced without possessing each chip. If extraction remains per-device and expensive, device
scarcity holds and the adversary must buy hardware. If one extraction generalises across a
family, the adversary obtains an oracle for minting attestations and the construction degrades
to precisely the identity farm it was built to prevent, having excluded honest modified devices
in the meantime.

This bears more heavily on the CA-free path than on the WebAuthn one, because that path has no
human tap to slow an oracle down. We are not aware of a satisfying public treatment at the
granularity the design requires, and we regard establishing it as the most valuable next step.
We would rather learn that the assumption fails than deploy on it silently.

\section{Related work}
\label{sec:related}

\textbf{Sybil resistance.} Douceur~\cite{douceur} establishes that without a trusted identity
authority, an adversary can present arbitrarily many identities absent a resource
constraint---and PoSEA is best read as choosing which resource. Social-graph
defences~\cite{sybilguard} bound identities by trust-graph structure, requiring a graph the
system must already possess. Proof-of-personhood~\cite{pop,offline-pseudonyms} binds identities
to physical gatherings; PoSEA binds them to physical objects, which scales without ceremony and
proves correspondingly less.

\textbf{Attestation without an issuer.} Direct Anonymous Attestation~\cite{daa} solves an
adjacent problem---removing the privacy CA from TPM attestation---by group signatures with an
issuer that sets up the group. It provides unlinkability, which we do not, and requires an
issuer and its setup, which we do not. Our construction removes the issuing party entirely, at
the cost of interaction and of a public ordering; it is applicable precisely because a ledger
supplies that ordering for free, and would not be applicable without one.

\textbf{Remote attestation generally.} SGX remote attestation~\cite{sgx} similarly terminates
in a vendor service, and shares the failure mode we report: when the vendor service is
unavailable to a party, the hardware's evidence becomes unusable regardless of its integrity.

\textbf{Commit--reveal.} The commitment scheme is textbook~\cite{blum}; the contribution is not
the scheme but Observation~\ref{obs:det}, that a TPM credential blob is deterministic in its
inputs, which is what makes a commit--reveal sufficient here. Without that, a challenge
could not be replayed and every verifier would need the challenger's word.

\section{Conclusion}

Per-identity admission rules cannot solve identity farming, because their cost is linear in
identities and therefore neutral between an honest participant and an adversary. The
alternative is to change what an identity is. Anchoring it in a physically distinct secure
element makes identity scarce rather than expensive.

For the quarter of machines whose platform will not certify their hardware, we show that the
certificate authority the standard construction requires can be removed entirely, because a TPM
credential blob is deterministic in its inputs and a ledger can supply ordering where an issuer
supplied a signature. The security of that construction rests on publication order and on one
object attribute, both of which a verifier checks itself.

The price is admitting vendor attestation authorities into the trust path, excluding owners who
have taken control of their hardware, and---on the CA-free path---resting on hardware price
rather than human effort. Whether that price is worth paying depends on an empirical question
about secure-element extraction costs that remains, to our knowledge, open.

\section*{Availability}

A standalone reference implementation, the specification, the pinned roots and the tests are
public at \url{https://github.com/hclivess/posea}; the enrolment of \S\ref{sec:noca} is
\texttt{posea/tpm.py} and \texttt{posea/enrolment.py}, specified normatively in
\texttt{SPEC.md}~\S9. The deployment it is extracted from is at
\url{https://github.com/hclivess/nado}. The reverse-engineered platform key-attestation claim
format is documented at \url{https://github.com/hclivess/windows-tpm-claim-format}. The network
is pre-mainnet; there is no token sale and nothing offered for purchase.


\begin{thebibliography}{99}

\bibitem{nakamoto}
S.~Nakamoto, ``Bitcoin: A peer-to-peer electronic cash system,'' 2008.

\bibitem{douceur}
J.~R. Douceur, ``The Sybil attack,'' in \emph{Proc. 1st Int. Workshop on Peer-to-Peer Systems
(IPTPS)}, 2002, pp.~251--260.

\bibitem{sybilguard}
H.~Yu, M.~Kaminsky, P.~B. Gibbons, and A.~Flaxman, ``SybilGuard: Defending against Sybil
attacks via social networks,'' in \emph{Proc. ACM SIGCOMM}, 2006.

\bibitem{pop}
M.~Borge, E.~Kokoris-Kogias, P.~Jovanovic, L.~Gasser, N.~Gailly, and B.~Ford,
``Proof-of-personhood: Redemocratizing permissionless cryptocurrencies,'' in \emph{IEEE
European Symp. on Security and Privacy Workshops (EuroS\&PW)}, 2017.

\bibitem{offline-pseudonyms}
B.~Ford and J.~Strauss, ``An offline foundation for online accountable pseudonyms,'' in
\emph{Proc. 1st Workshop on Social Network Systems (SocialNets)}, 2008.

\bibitem{casper}
V.~Buterin and V.~Griffith, ``Casper the friendly finality gadget,'' arXiv:1710.09437, 2017.

\bibitem{idena}
Idena, ``Idena: Proof-of-person blockchain.'' [Online]. Available: \url{https://idena.io}

\bibitem{webauthn}
W3C, ``Web Authentication: An API for accessing public key credentials,'' W3C Recommendation.
[Online]. Available: \url{https://www.w3.org/TR/webauthn-3/}

\bibitem{androidkeyattest}
Google, ``Verifying hardware-backed key pairs with key attestation,'' Android Developers.
[Online]. Available:
\url{https://developer.android.com/privacy-and-security/security-key-attestation}

\bibitem{tpm2}
Trusted Computing Group, ``Trusted Platform Module Library, Part 1: Architecture'' and
``Part 3: Commands,'' Family 2.0, Level 00, Rev. 01.59, 2019.

\bibitem{tcgek}
Trusted Computing Group, ``TCG EK Credential Profile for TPM Family 2.0,'' Level 0, Version
2.5.

\bibitem{daa}
E.~Brickell, J.~Camenisch, and L.~Chen, ``Direct anonymous attestation,'' in \emph{Proc. 11th
ACM Conf. on Computer and Communications Security (CCS)}, 2004, pp.~132--145.

\bibitem{sgx}
V.~Costan and S.~Devadas, ``Intel SGX explained,'' IACR Cryptology ePrint Archive, Report
2016/086, 2016.

\bibitem{blum}
M.~Blum, ``Coin flipping by telephone: A protocol for solving impossible problems,'' in
\emph{Proc. CRYPTO}, 1981, pp.~11--15.

\bibitem{rfc8017}
K.~Moriarty, B.~Kaliski, J.~Jonsson, and A.~Rusch, ``PKCS \#1: RSA cryptography
specifications version 2.2,'' RFC 8017, IETF, 2016.

\bibitem{finney}
H.~Finney, ``Bleichenbacher's RSA signature forgery based on implementation error,''
openssl-dev mailing list, Aug. 2006.

\bibitem{cve2006}
``CVE-2006-4339: OpenSSL RSA signature forgery,'' 2006.

\bibitem{tpmfail}
D.~Moghimi, B.~Sunar, T.~Eisenbarth, and N.~Heninger, ``TPM-FAIL: TPM meets timing and lattice
attacks,'' in \emph{Proc. 29th USENIX Security Symp.}, 2020.

\bibitem{faultpm}
H.~N. Jacob, C.~Werling, R.~Buhren, and J.-P. Seifert, ``faulTPM: Exposing AMD firmware TPMs,''
arXiv:2304.14717, 2023.

\bibitem{trustzone}
A.~Shakevsky, E.~Ronen, and A.~Wool, ``Trust dies in darkness: Shedding light on Samsung's
TrustZone keymaster design,'' in \emph{Proc. 31st USENIX Security Symp.}, 2022.

\bibitem{claimformat}
J.~Ku\v{c}era, ``Windows TPM key-attestation claim format,'' 2026. [Online]. Available:
\url{https://github.com/hclivess/windows-tpm-claim-format}

\end{thebibliography}

\end{document}
